\chapter{多分类学习的可学习性}
\label{chap:multiclass-learnability}

一个二分类假设改答时，只剩另一种标签可选；多分类却可能改成多个不同答案。候选在有限输入上的组合因此增加，前章将可学习性与一致收敛合在一起的刻画也需要重查。标签有限时，我们先扩展预测计数；允许无限标签后，分别判断全类风险估计与成功学习器是否存在；再为统计结果增加程序条件。

本章把第~\ref{sec:learning-problem}节的二类标签改为有限或无限$Y$，仍观察输入标签对，使用i.i.d.抽样、零一风险和前章PAC基准。全类损失的共同估计由\term{图维}（graph dimension）描述，成功学习器的存在由\term{DS维}（Daniely--Shalev-Shwartz dimension）描述。一般多分类允许输出空间$\mathcal W_{\mathrm{out}}$超出比较类$\mathcal H$，这是后述学习刻画的条件。

\section{从二分类到多分类}
\label{sec:multiclass-setting}

邮件标签现在分正常、广告、诈骗，输入$x$仍是特征，$h(x)$选择类别；与真实$y$不同即计一次错误。以$X,Y$为两空间、非空$\mathcal H\subseteq Y^X$为候选，对联合来源$\mathcal P$及iid样本$S=((x_i,y_i))_{i=1}^n$，定义~\ref{def:general-population-risk}与定义~\ref{def:general-empirical-risk}在零一损失下为
\begin{equation}
 R_{\mathcal P}(h)=\mathbb P_{(x,y)\sim\mathcal P}\{h(x)\neq y\},\qquad
 \hat R_S(h)=\frac1n\sum_{i=1}^n\mathbb I\{h(x_i)\neq y_i\}.
 \label{eq:multiclass-risk}
\end{equation}
后文固定这一计分口径。可实现PAC要求零风险目标来源下输出风险小，不可知PAC要求接近类内下确界；量词沿用第~\ref{sec:pac-learning}节。统计定理仍须满足第~\ref{sec:learning-problem}节的可测性与可分性条件。

有限标签可先编码成二元位。例如对三类取
\[
 \text{正常}\longleftrightarrow(0,1),\qquad
 \text{广告}\longleftrightarrow(1,0),\qquad
 \text{诈骗}\longleftrightarrow(1,1).
\]
广告邮件于是给两位任务标签1和0。在同一批邮件上训练$b_1,b_2$，新输入先得到$(b_1(x),b_2(x))$，再按表解码；无效$(0,0)$预定解为正常。这样改成两个二类任务，类别信息并未减少。

当$2\leq K=|Y|<\infty$时，用$L=\lceil\log_2K\rceil$位赋互异编码$c:Y\to\{0,1\}^L$，以$c_j(y)$记第$j$位。相应二类候选必须来自原空间的这一位投影：
\[
 \mathcal B_j=\{x\mapsto c_j(h(x)):h\in\mathcal H\},\qquad 1\leq j\leq L.
\]
每个$h$先编码输出，再保留第$j$位。标签虽二元，诱导类仍可复杂；二类分布无关PAC须有有限VC，编码本身未保证这一容量条件。若$K=1$，唯一标签的常量预测已是零风险规则，无需拆分位任务。

可实现来源由$h^\star\in\mathcal H$给零风险时，$c_j\circ h^\star$在每个位类中且几乎处处预测$c_j(y)$，所以全部位任务也可实现。若每个$\mathcal B_j$可PAC学习，令学得$b_j$的位风险为
\[
 e_j=\mathbb P_{(x,y)\sim\mathcal P}\{b_j(x)\neq c_j(y)\}.
\]
以$\tilde h$记解码分类器。全部位正确必给正确类别，故类别错误包含在位错并集中：
\[
 R_{\mathcal P}(\tilde h)\leq\sum_{j=1}^L e_j.
\]
将位风险目标定为$\varepsilon/L$、位学习失败预算定为$\delta/L$，两次并集分别得到类别风险与抽样成功概率，从而以至少$1-\delta$的概率使风险至多$\varepsilon$。位学习可共享样本，错误不必独立，只需共同样本量满足
\[
 n\geq\max_{1\leq j\leq L}
 m_{\mathcal B_j}(\varepsilon/L,\delta/L).
\]
合成输出未必属于原$\mathcal H$。

逐位学习在不可知情形不一定优化原任务。沿用三类编码，设输入恒定，正常、广告、诈骗的概率为0.40、0.35、0.25，候选为三个常量假设。第一位为1的概率0.60，第二位为1的概率0.65，各位最优都选1，合成诈骗后风险0.75；直接选正常的风险却为0.60。即使已知所有概率，逐位最优仍选错原任务的最优类别。逐位优化分别降低每个编码位的错误概率；多类零一损失只在整个标签正确时计0，其余一律计1。同一错误标签可能在两位上都出错，两项位损失便分别记错，而类别损失只记一次。两个优化目标因此并不相同。

有限标签编码给出算法构造，学习保证还须控制诱导二类空间并合并位风险。无限标签无法由任何固定有限位数区分，上述有限次归约也不能直接使用。要直接判断原空间，我们转而分析它在有限输入上产生的预测结构。

这里不能把“实现全部$|Y|^m$种预测”作为打散要求。有些空间在各输入只用两个标签，却已包含任意大的二类难题，要求用齐所有标签反而漏掉它。我们应先固定各输入的两种答案，再问能否独立切换；\term{Natarajan维}（Natarajan dimension）衡量可以维持这种自由选择的输入数。

\section{有限标签空间及Natarajan维的边界}
\label{sec:multiclass-dimensions}

先限定$2\leq K=|Y|<\infty$。找出独立二元选择，再用其最大规模限制预测数，便能沿用前章的对称化与ERM风险比较。

\subsection{Natarajan维}

固定输入处的两种标签选择能够限制有限标签类的预测数。Natarajan打散要求先为每个输入固定两个不同答案，再由类中的函数实现这组标签对的全部独立选择；标签对可以随输入改变，却不能随待实现的模式改变（可参考\BookDefinitionRef{def:natarajan-dimension}{第一卷“计数与计算”章的Natarajan打散与维度定义}）。本章以$d_N(\mathcal H)$记可打散有限输入集大小的上确界。
% 旧公式位置保留为导航锚点，不再重复通用定义式。
\phantomsection\label{eq:learning-natarajan-shattering}

这一条件不要求标签有限。二类时各处的不同标签对只能取那两个答案，故它退化为VC维。维度0说明各输入预测唯一；可打散点集没有有限规模上界时，维度为无穷。这里继续把所得维度接到预测计数。

\subsection{有限标签的一致收敛}

固定输入后先合并相同预测。令
\[
 V(x_1,\ldots,x_m)
 =\{(h(x_1),\ldots,h(x_m)):h\in\mathcal H\}.
\]
再对输入选择取上确界，得到生长函数
\[
 \Pi_{\mathcal H}(m)=\sup_{(x_1,\ldots,x_m)\in X^m}|V(x_1,\ldots,x_m)|.
\]
这与第~\ref{sec:binary-growth-sauer}节的二类计数步骤相同，区别是每坐标可取$K$个标签。Sauer--Shelah由VC限制二元数目，下面用Natarajan承担多类的对应职责。

% 保留旧位置锚点；这里直接使用组合计数结论。
\phantomsection\label{thm:learning-multiclass-growth-bound}
当$2\leq K=|Y|<\infty$且$d=d_N(\mathcal H)<\infty$时，记$q=K(K-1)/2$为不同无序标签对的数目。对每个整数$m\geq0$，有限标签的计数界给出
\[
 \Pi_{\mathcal H}(m)\leq G_{d,K}(m)
 :=\sum_{j=0}^{\min\{d,m\}}\binom mj q^j
\]
（可参考\BookTheoremRef{thm:multiclass-growth-bound}{第一卷“计数与计算”章的有限标签预测数目上界}）。在学习问题中，$G_{d,K}(2n)$将控制原样本与副样本合并后的不同损失模式数，因而能够进入下一步对称化后的失败概率。

二类时$q=1$、$d_N=\operatorname{VC}(\mathcal H)$，求和回到Sauer--Shelah的$\sum_{j=0}^{\min\{d,m\}}\binom mj$。多类的$q^j$另计每个参与二元选择的坐标所用的标签对。

固定$d,K$使上界至多多项式增长。样本标签一经固定，相同预测向量也给出相同损失向量，故这个计数控制经验风险分析中的不同结果。副样本与随机交换于是可以沿用。

% 避免短定理在页末仅留下标题和半句条件。
\begingroup
\tcbset{lines before break=4}
\begin{theorem}[有限标签的一致收敛保证]
\label{thm:finite-label-uniform-convergence}
设$(X,\mathcal X)$为可测空间，离散标签空间$Y$满足
$2\leq|Y|=K<\infty$，$\mathcal H\subseteq Y^X$是由可测分类器组成的非空集合，
$d=d_N(\mathcal H)<\infty$，并令
\[
 G_{d,K}(m)=\sum_{j=0}^{\min\{d,m\}}\binom mj
 \left(\frac{K(K-1)}2\right)^j.
\]
假定零一损失类$\mathcal L_{\mathcal H}
=\{(x,y)\mapsto\mathbb I\{h(x)\neq y\}:h\in\mathcal H\}$中的每个函数均可测，
并存在可数子类$\mathcal L_0\subseteq\mathcal L_{\mathcal H}$，使每个
$f\in\mathcal L_{\mathcal H}$都是$\mathcal L_0$中某个函数序列的逐点极限。
设$\mathcal P$为$X\times Y$上的任意联合分布，$n\geq1$，
$S=((x_i,y_i))_{i=1}^n\sim\mathcal P^n$，并记
\[
 R_{\mathcal P}(h)=\mathbb P_{(x,y)\sim\mathcal P}\{h(x)\neq y\},\qquad
 \hat R_S(h)=\frac1n\sum_{i=1}^n\mathbb I\{h(x_i)\neq y_i\}.
\]
则对任意$\delta\in(0,1)$，以关于$S\sim\mathcal P^n$至少$1-\delta$的概率，
\begin{equation}
 \sup_{h\in\mathcal H}|R_{\mathcal P}(h)-\hat R_S(h)|
 \leq\sqrt{\frac8n\log\frac{4G_{d,K}(2n)}{\delta}}.
 \label{eq:finite-label-uniform-bound}
\end{equation}
\end{theorem}
\endgroup

\begin{proof}
取$\eta>0,n\eta^2\geq2$。独立副样本把总体/经验偏差转成双经验风险差超过$\eta/2$，概率最多增一倍。合并输入上的损失向量至多$G_{d,K}(2n)$种；逐对随机交换后，对固定向量集中，再对这些有限向量求并集界，得到
\[
 \mathbb P\!\left\{\sup_h|R_{\mathcal P}(h)-\hat R_S(h)|>\eta\right\}
 \leq4G_{d,K}(2n)e^{-n\eta^2/8}.
\]
令右侧等于$\delta$即可解出所列偏差；该解自动满足$n\eta^2\geq2$。上界超过1时，零一损失的取值范围已保证结论。
\end{proof}

这里$\mathcal P$仍是联合来源，没有要求各类别等概率。固定$d,K,\delta$后，式~\eqref{eq:finite-label-uniform-bound}随$n$趋零且统一覆盖全部允许来源，因此得到分布无关UC。当$1\leq d\leq2n$时，代入$G_{d,K}(2n)\leq q^d(2en/d)^d$还能看出类别数以$\log K$量级进入上界。

\subsection{有限标签的可学习性}

共同风险估计可以覆盖样本选出的ERM。用输出与参照各一次偏差，再结合Natarajan打散产生的困难来源，便得到有限标签下的双向刻画。

\begin{theorem}[有限标签的可学习性刻画]
\label{thm:finite-label-learnability}
设$(X,\mathcal X)$为可测空间，离散标签空间$Y$满足
$2\leq|Y|=K<\infty$，$\mathcal H\subseteq Y^X$是由可测分类器组成的非空集合，
并采用零一损失。假定损失类
$\mathcal L_{\mathcal H}
=\{(x,y)\mapsto\mathbb I\{h(x)\neq y\}:h\in\mathcal H\}$中的每个函数均可测，
且存在可数子类$\mathcal L_0\subseteq\mathcal L_{\mathcal H}$，使每个
$f\in\mathcal L_{\mathcal H}$都是$\mathcal L_0$中某个函数序列的逐点极限。
学习算法均为可测映射，训练样本均为$S\sim\mathcal P^n$，成功概率同时计入样本抽取
和算法的内部随机性。在这一分布无关设定中，以下命题等价：
\begin{enumerate}
 \item $d_N(\mathcal H)<\infty$；
 \item 零一损失满足分布无关的一致收敛；
 \item $\mathcal H$不可知PAC可学习；
 \item $\mathcal H$可实现PAC可学习。
\end{enumerate}

设$d=d_N(\mathcal H)<\infty$，并令
\[
 G_{d,K}(m)=\sum_{j=0}^{\min\{d,m\}}\binom mj
 \left(\frac{K(K-1)}2\right)^j.
\]
对任意$\varepsilon,\delta\in(0,1)$、任意$X\times Y$上的联合分布$\mathcal P$，
若$n\geq1$且
\[
 n\geq\frac{32}{\varepsilon^2}\log\frac{4G_{d,K}(2n)}{\delta},
\]
则任意以概率1输出
$A(S)\in\argmin_{h\in\mathcal H}\hat R_S(h)$的可测ERM，在
$S\sim\mathcal P^n$及其内部随机性的联合概率下，以至少$1-\delta$的概率满足
\[
 R_{\mathcal P}(A(S))\leq\inf_{h\in\mathcal H}R_{\mathcal P}(h)+\varepsilon.
\]
\end{theorem}

\begin{proof}
有限Natarajan维先给UC。取偏差$\varepsilon/2$，在共同成功事件内用ERM最优性比较任意$h$：
\begin{align*}
 R_{\mathcal P}(A(S))&\leq\hat R_S(A(S))+\varepsilon/2\\
 &\leq\hat R_S(h)+\varepsilon/2\\
 &\leq R_{\mathcal P}(h)+\varepsilon.
\end{align*}
对参照取下确界得不可知PAC。可实现分布属于这一允许族，因此也得到可实现PAC。

逆向设Natarajan无限。任意$n$都可取被打散的$2n$点及各点两种固定标签，令输入均匀，目标在各标签对中独立公平选择。$n$次观测至多见$n$点，余下标签仍公平，所以随机目标下任意算法平均风险至少$1/4$。平均中必有固定可实现目标也达此值；再由$R\leq1$及$\mathbb E R\leq1/8+(7/8)\mathbb P\{R>1/8\}$，得失败概率至少$1/7$。每个$n$都存在这种来源，违反分布无关可实现PAC。
\end{proof}

必要性证明只使用各输入固定的两个标签，没有用标签总数有限，所以任意$Y$下可学习都必须有有限Natarajan维。有限标签承担的是另一方向的预测计数。ERM样本式是充分条件，尚未声称匹配两类PAC的最优复杂度。

\subsection{Natarajan维在无限标签空间下的局限}
\label{sec:multiclass-erm-limitations}

无限$Y$使$q=\binom K2$不再有限，旧计数界失去作用。证明失效本身不等于维度判据失效；还须构造有限Natarajan维却不可学习的空间，才能否定充分性。

先在两坐标看一个有限模式：
\[
 (a,u),\ (b,u),\ (b,v),\ (c,v),\ (c,w),\ (a,w),
\]
其中$a,b,c$互异，$u,v,w$互异。每个向量在每坐标都有只改该坐标的邻居，例如$(a,u)$的$(b,u)$和$(a,w)$。但任意两种第一标签至多共享一种第二标签，无法形成固定两对标签的四组合；Natarajan维仅1，两坐标却都留歧义。这一个有限类仍可学习，反例还需任意维的同类结构。

图~\ref{fig:learning-theory-ds-neighbors}把这六种预测的邻接关系画成环。每个顶点都连接一条只改第一坐标的边和一条只改第二坐标的边，因此两坐标上的局部歧义可以同时保留；这种邻接条件不要求固定两对标签的四种组合全部出现。

\begin{figure}[htbp]
  \centering
  \includegraphics[width=93.405mm]{figures/02-foundations/02-learning-theory/tikz/learning-theory-ds-neighbors/figure.pdf}
  \caption{六种预测的坐标邻接环。顶点为上文的六个向量，蓝色实线仅改变第一坐标，绿色虚线仅改变第二坐标；边没有方向。每个顶点在两个坐标上分别有邻居，却不存在固定两对标签的完整四种组合。因而这个两输入函数类的Natarajan维为1，后文分析的DS维为2；它仍是可学习的有限类，不能单独证明不可学习性。}
  \label{fig:learning-theory-ds-neighbors}
\end{figure}

若要从二维局部歧义得到不可学习性，需要任意维的对应构造。下面把$V_m$的向量视为坐标集$\{1,\ldots,m\}$上的函数，使其Natarajan维有明确对象。

\begin{lemma}[低Natarajan维的局部邻接构造]
\label{lem:low-natarajan-pseudocubes}
对每个整数$m\geq1$，把$Y_m^m$中的向量视为定义在坐标集
$\{1,\ldots,m\}$上的分类函数。存在非空有限标签集$Y_m$及有限非空函数集合
$V_m\subseteq Y_m^{\{1,\ldots,m\}}$，使
\begin{enumerate}
 \item $d_N(V_m)=1$；
 \item 对任意$v\in V_m$和$i\in\{1,\ldots,m\}$，存在$u\in V_m$满足
 \[
 u_i\neq v_i,\qquad u_j=v_j\quad(j\neq i).
 \]
\end{enumerate}
\end{lemma}

Brukhim等的组合拓扑构造给出该引理\scite{brukhim2022multiclass}。证明未在本章展开。它同时限制固定标签对的自由选择，又保留每个向量在每个坐标的竞争答案。依靠这两项性质，下面把各维局部构造合成一个固定假设空间。

\begin{theorem}[无限标签下有限Natarajan维并不充分]
\label{thm:natarajan-infinite-counterexample}
在输入与标签均取离散$\sigma$代数、训练样本独立同分布并采用零一损失的设定中，
存在可数输入空间$X$、可数标签空间$Y$和非空假设空间
$\mathcal H\subseteq Y^X$，使$d_N(\mathcal H)=1$，但不存在对所有
$\mathcal H$可实现联合分布都具有分布无关可实现PAC保证的可测学习算法。
\end{theorem}

\begin{proof}
对各$m$取引理~\ref{lem:low-natarajan-pseudocubes}中的$V_m$，把输入组成互不相交的块$C_m=\{(m,1),\ldots,(m,m)\}$，令$X=\bigcup_{m\geq1}C_m$。局部标签加块编号写为$(m,y)$，因此$(2,a)$与$(3,a)$是不同类别。另设普通标签$\star$表示不在所选块，它不同于全部$(m,y)$，并非未知、缺失或任意值。标签总空间为
\[
 Y=\{\star\}\cup\bigcup_{m\geq1}(\{m\}\times Y_m).
\]
每个$v\in V_m$扩展为整个输入空间上的假设
\[
 h_{m,v}(r,i)=
 \begin{cases}
 (m,v_i),&r=m,\\
 \star,&r\neq m.
 \end{cases}
\]
例如$h_{2,v}$只在$C_2$使用局部答案，其他块均答$\star$。以全部$h_{m,v}$组成$\mathcal H$，各维困难块便属于同一个固定空间。

两点若在不同块，每处标签对至少含一个非$\star$标签，而任一假设只在一个块输出这种标签，故无法实现两处同时非$\star$的组合。两点若在同块，标签对含$\star$时，一处为$\star$会使该块全为$\star$，不能独立切换；两对都不含$\star$时，又要求$V_m$打散两坐标，与维度1矛盾。单点仍可取局部标签或$\star$，所以$d_N(\mathcal H)=1$。

对给定$n$取$m=2n$，输入在$C_m$均匀，目标从有限$V_m$均匀抽取。考虑未见坐标$i$，即使透露其他全部坐标的真实标签，该处仍有竞争答案。给定其余坐标后，不同向量在$i$处必有不同标签；条件分布对这些向量均匀，邻接条件又保证至少两个。因此该处的平均错误概率至少$1/2$。

$n$次抽样至少遗漏$2n$点的一半，故对目标、样本及内部随机性平均的风险至少$1/4$。有限目标平均中必有固定$v$达到该下界；用$\mathbb E R\leq1/8+(7/8)\mathbb P\{R>1/8\}$，再得风险超过$1/8$的概率至少$1/7$。每个样本量都有如此可实现来源，所以不存在统一PAC样本量。
\end{proof}

无限标签本身没有排除学习；障碍是Natarajan有限未排除所有困难结构。接下来把全类共同估计与成功算法存在分开，寻找各自的容量条件。

\section{图维、DS维与一般标签空间}
\label{sec:general-multiclass-dimensions}
\label{sec:multiclass-sample-complexity}
有限标签靠Natarajan的预测计数连接UC与学习。去掉标签数有限的条件后，第~\ref{sec:multiclass-erm-limitations}节已给维度有限却不可学习的例子。我们因此分别问：一份样本能否同时估准所有风险，以及是否存在算法选出低风险者？前者约束整个空间，后者只约束算法输出。

零一风险只问是否匹配标签，图维因此可转为二值损失类的VC问题。学习器则可能利用标签携带的更多信息，DS维用预测之间单坐标变化描述剩余歧义。先建立图维与UC的关系，再分析为什么存在学习器是另一项要求。

\subsection{一致收敛}

零一损失把无限个错误答案都计为1，风险只看与真值是否匹配。为描述这一二元能力，我们先给每个输入固定参考答案$f(x)$，再问能否任意指定一组匹配位置。\term{图打散}（graph shattering）要求对每个位置子集都有类中函数恰在这些位置匹配参考，在其余位置全部偏离参考；偏离时使用的标签可以随模式改变（可参考\BookDefinitionRef{def:graph-dimension}{第一卷“计数与计算”章的图打散定义}）。

可图打散有限输入集大小的上确界记为$d_G(\mathcal H)$（可参考\BookDefinitionRef{def:graph-dimension-value}{第一卷“计数与计算”章的图维定义}）。图打散只预定参考答案，Natarajan打散还预定每处的替代答案，因此后者推出前者，逆向未必成立。接下来利用“匹配或偏离”的二元模式，将图维转为零一损失类的VC维。

风险要看匹配与否，故将各假设改写为$\mathcal L_{\mathcal H}=\{(x,y)\mapsto\mathbb I\{h(x)\neq y\}:h\in\mathcal H\}$。新函数以输入标签对为自变量，以0或1为输出。

\begin{theorem}[图维与一致收敛]
\label{thm:graph-uniform-convergence}
设$(X,\mathcal X)$为可测空间，$Y$为取离散$\sigma$代数的非空标签空间，
$\mathcal H\subseteq Y^X$是由可测分类器组成的非空集合，并令
\[
 \mathcal L_{\mathcal H}
 =\{(x,y)\mapsto\mathbb I\{h(x)\neq y\}:h\in\mathcal H\}.
\]
则$\operatorname{VC}(\mathcal L_{\mathcal H})=d_G(\mathcal H)$。
进一步假定$\mathcal L_{\mathcal H}$中的函数均可测，且存在可数子类
$\mathcal L_0\subseteq\mathcal L_{\mathcal H}$，使每个损失函数都是
$\mathcal L_0$中某个函数序列的逐点极限。此时$d_G(\mathcal H)<\infty$当且仅当：
对任意$\eta,\delta\in(0,1)$，存在不依赖联合分布$\mathcal P$的样本量
$m(\eta,\delta)$，使任意$n\geq m(\eta,\delta)$及
$S=((x_i,y_i))_{i=1}^n\sim\mathcal P^n$均满足
\[
 \mathbb P_{S\sim\mathcal P^n}\left\{
 \sup_{h\in\mathcal H}|R_{\mathcal P}(h)-\hat R_S(h)|\leq\eta
 \right\}\geq1-\delta,
\]
其中$R_{\mathcal P}$与$\hat R_S$均为零一损失风险。
\end{theorem}

\begin{proof}
若以$f$图打散$C$，在$\{(x,f(x)):x\in C\}$上，匹配/不匹配使损失实现全部二元赋值，故$\operatorname{VC}(\mathcal L_{\mathcal H})\geq d_G$。

反过来，损失打散的点集不能含同一$x$配两种$y$，因为两处损失不可能同时为0。输入于是互异，置$f(x_i)=y_i$，损失0/1便表示匹配/偏离，各模式使这些输入图打散。两维相等后，前章二值函数类的VC/UC刻画即给出所述等价。
\end{proof}

图维现在判定全类风险是否可共同估准。共同事件也覆盖ERM所选输出，便可将经验最优转换为接近类内最低总体风险。下面列出这一推论，并核对两次偏差的来源。

\begin{theorem}[一致收敛推出可学习性]
\label{thm:multiclass-uniform-pac}
设$(X,\mathcal X)$为可测空间，$Y$为离散标签空间，
$\mathcal H\subseteq Y^X$是由可测分类器组成的非空集合。采用零一损失，记
\[
 R_{\mathcal P}(h)=\mathbb P_{(x,y)\sim\mathcal P}\{h(x)\neq y\},\qquad
 \hat R_S(h)=\frac1n\sum_{i=1}^n\mathbb I\{h(x_i)\neq y_i\}.
\]
假定损失类中的函数均可测，并存在可数损失子类，使每个损失函数都是该子类中某个
函数序列的逐点极限；$\mathcal H$满足分布无关的一致收敛。对任意
$\varepsilon,\delta\in(0,1)$，若$n\geq1$足以使任意$X\times Y$上的联合分布
$\mathcal P$及样本$S\sim\mathcal P^n$满足
\[
 \mathbb P\!\left\{
 \sup_{h\in\mathcal H}|\hat R_S(h)-R_{\mathcal P}(h)|\leq\varepsilon/2
 \right\}\geq1-\delta,
\]
则任意以概率1输出
$A(S)\in\argmin_{h\in\mathcal H}\hat R_S(h)$的可测ERM学习算法$A$，均以至少
$1-\delta$的概率满足
\[
 R_{\mathcal P}(A(S))\leq\inf_{h\in\mathcal H}R_{\mathcal P}(h)+\varepsilon.
\]
概率同时来自$S\sim\mathcal P^n$和$A$的内部随机性。
因此，$\mathcal H$不可知PAC可学习，也可实现PAC可学习。
\end{theorem}

\begin{proof}
将每个$h$换成二值损失函数，经验最优比较保持不变。第~\ref{chap:binary-learnability}章的两次$\varepsilon/2$偏差分别用于输出和任意参照，得到总体超额至多$\varepsilon$；再对参照取下确界即为不可知PAC。可实现来源使该下确界为0，因而也满足可实现保证。
\end{proof}

图维有限经UC再经ERM，足以推出学习；无限标签下，它却不再是学习存在的必要条件。

UC约束整个空间，学习器可能只需要辨明正确目标。下面让标签本身透露目标信息，展示全类不UC却有成功算法的情形。

这个标签传信息的构造取$X=\mathbb N$，用每个有限自然数子集作一个类别，再设与它们均不同的$\star$表示输入在目标集合外。对有限$A$定义
\begin{equation}
 h_A(x)=\begin{cases}A,&x\in A,\\\star,&x\notin A,\end{cases}
 \qquad\mathcal H_C=\{h_A:A\subset\mathbb N\text{有限}\}.
 \label{eq:cantor-class}
\end{equation}
标签$\{2,5\}$直接列明集合成员，所以在$\mathcal H_C$可实现来源下，一次非$\star$观测就能恢复目标；样本全为$\star$时输出$h_{\varnothing}$。这里的数学标签故意携带整条规则，与通常的语义类别不同。

\begin{theorem}[可学习性不推出一致收敛]
\label{thm:multiclass-erm-separation}
令$X=\mathbb N$，标签空间$Y$由特殊标签$\star$与$\mathbb N$的所有有限子集组成，
二者均取离散$\sigma$代数。对每个有限$A\subseteq\mathbb N$定义
\[
 h_A(x)=\begin{cases}A,&x\in A,\\\star,&x\notin A,\end{cases}
 \qquad
 \mathcal H_C=\{h_A:A\subseteq\mathbb N\text{有限}\}.
\]
采用独立同分布抽样与零一损失。学习算法在样本出现非$\star$标签$A$时输出$h_A$，
在所有标签均为$\star$时输出$h_{\varnothing}$；该算法使$\mathcal H_C$可实现PAC可学习。
另一方面，对大小为$n$且标签全为$\star$的样本，令
$A_S=\{1,\ldots,2n\}\setminus\{x_1,\ldots,x_n\}$并输出$h_{A_S}$；对其他样本，
在$h_{\varnothing}$及样本中出现的有限集合标签对应的假设之间，按预先固定的顺序选择
零一经验风险最小者。这个规则是确定性恰当ERM，却没有分布无关PAC保证。
此外，$\mathcal H_C$满足$d_N(\mathcal H_C)=1$及$d_G(\mathcal H_C)=\infty$，
因而不满足零一损失的一致收敛。
\end{theorem}

\begin{proof}
设目标$h_A$，以$p=\mathcal P_X(A)$记其集合概率。见过非$\star$时恢复目标，风险0；否则输出空集合规则，风险$p$。若$p>\varepsilon$，遗漏事件概率$(1-p)^n\leq e^{-n\varepsilon}$，所以$n\geq\log(1/\delta)/\varepsilon$足够。

任意有限$C$以全$\star$为参考都可图打散：要恰在$B$输出$\star$，取$A=C\setminus B$。图维因而无限，UC失败。非$\star$标签却唯一确定$h_A$；若两点被Natarajan打散，固定一处标签对中的非$\star$者就固定整个假设，另一处无法再自由选两标签，矛盾。单点可取两个标签，所以$d_N=1$。

这个空间仍有有限样本ERM。若集合标签$A$未在样本出现，$h_A$相比$h_{\varnothing}$不会修正任何错误，却可能把正确的$\star$改错。因此经验最小值可在$h_{\varnothing}$与样本出现的有限个集合标签对应假设中取得。预先为全部有限集合排一个顺序，即可固定平局选择。

让另一ERM在$n$条全$\star$记录时故意选$h_{A_S}$，其中$A_S=\{1,\ldots,2n\}\setminus\{x_1,\ldots,x_n\}$；其他可实现样本从集合标签恢复目标，不可实现样本则按固定顺序取前述有限候选的经验最小者。全$\star$样本上它零经验错，但目标为空、输入均匀于$\{1,\ldots,2n\}$时，至少$n$个未见输入被误判，风险至少$1/2$。任意$n$都有这样的来源，这个ERM不能分布无关PAC学习。
\end{proof}

好坏两算法在可实现样本上都达到零经验风险，区别在于怎样使用标签信息。这里“有成功学习器”不能由“所有ERM都成功”代替\scite{daniely2015multiclass}。第~\ref{sec:multiclass-erm-limitations}节的空间完全没有PAC学习器，这个空间则只是某些经验最优选择失败，两例承担不同边界。

\subsection{DS维与可学习性}
\label{sec:ds-dimension}

集合标签例说明图维无限仍可学习，另一例的Natarajan维虽有限，却没有PAC学习器。所需充要结构应比固定标签对宽，又不要求全类风险估准。六预测例的线索是：即使知道其他输入的答案，遗漏位置也可能仍有竞争标签。

为了描述遗漏坐标上的竞争答案，我们在限制类中寻找有限非空的预测集合$V\subseteq Y^m$。要求每个$v\in V$在每个坐标$i$上都有集合内的另一个向量$u$：它只改变该坐标的答案，其余全部坐标与$v$相同。这正是本章下界证明需要的\term{伪立方体}（pseudo-cube）条件（可参考\BookDefinitionRef{def:pseudocube}{第一卷“计数与计算”章的伪立方体定义}）。

六种预测给二维伪立方体，全部二元$2^m$预测给$m$维伪立方体。多类不要求各坐标固定两种标签，却仍要求预测集合有限；无限邻接集的性质可能不同，不能删除这个条件。

DS维寻找能保留上述逐坐标歧义的最大输入规模。对$m$个互异输入按固定次序组成的$C$，检查限制类$\mathcal H|_C$是否含$m$维伪立方体；满足条件的有限输入规模取上确界，记为$d_{DS}(\mathcal H)$（可参考\BookDefinitionRef{def:ds-dimension}{第一卷“计数与计算”章的DS维定义}）。下面考察它如何决定学习。

图~\ref{fig:learning-theory-ds-neighbors}展示第~\ref{sec:multiclass-erm-limitations}节的六种预测，两种线型分别表示只改变一个坐标的邻接。每个顶点在两个坐标上均有邻居，所以这个两输入类的DS维为2；环缺少固定两标签对的完整四种组合，所以Natarajan维仅1。这个有限类仍可学习；无限标签下的不可学习反例还依赖任意维的局部构造及其在同一假设空间中的合并。

在$m-1$个坐标已知时，剩余邻居使遗漏标签未定。对每个向量、每个坐标都要求邻居，避免仅在个别向量周围留少量分支。固定标签对、固定参考答案与逐坐标邻接因而是不同要求；图中的一段邻接关系只有扩展到有限集合内的每个向量和每个坐标，才满足完整伪立方体条件（可参考\BookSectionRef{sec:comb-multiclass-dimensions}{第一卷“计数与计算”章的“多值函数族的容量”部分}）。

% 保留旧位置锚点；通用维度比较引用卷一。
\phantomsection\label{thm:learning-multiclass-dimension-comparison}
这些结构满足$d_N(\mathcal H)\leq d_{DS}(\mathcal H)\leq d_G(\mathcal H)$；二标签时三者均等于VC维。有限标签空间下，Natarajan维有限还蕴含图维有限（可参考\BookTheoremRef{thm:multiclass-dimension-comparison}{第一卷“计数与计算”章的多分类维度比较定理}）。这里使用三个关系：固定标签对的全部组合构成伪立方体；逐坐标邻接支持图打散；有限标签的指数图模式又受多项式预测计数限制。

有限标签的更精细界为$d_G=O(d_N\log K)$，不能外推无限$Y$。第~\ref{sec:multiclass-erm-limitations}节的不可学习反例有$d_N=1,d_{DS}=\infty$，故第一道维度差可以无穷。集合标签类$\mathcal H_C$则有$d_{DS}=1$：出现非$\star$的向量唯一确定假设，不能在别的坐标留邻居，全$\star$独自也不是伪立方体。结合单点可打散得维度1，而图维无限，故第二道也不能反向替换。

任意坐标都留竞争答案，使未见位置无法由其他标签唯一确定。若任意大的有限输入集都具有这项结构，就能为每个样本量构造学习失败的概率下界。

\begin{theorem}[无限DS维的不可学习性]
\label{thm:infinite-ds-lower-bound}
设$(X,\mathcal X)$为可测空间，$Y$为离散标签空间，
$\mathcal H\subseteq Y^X$是由可测分类器组成的非空集合，并采用零一损失。
若$d_{DS}(\mathcal H)=\infty$，则对每个整数$n\geq1$及任意可测学习算法$A$
（允许内部随机性，并允许输出$\mathcal H$之外的可测分类器），都存在
$X\times Y$上的联合分布$\mathcal P$及$h^\star\in\mathcal H$，使
$R_{\mathcal P}(h^\star)=0$，且对$S\sim\mathcal P^n$有
\[
 \mathbb P_{S,A}\{R_{\mathcal P}(A(S))>1/8\}\geq1/7.
\]
这里$R_{\mathcal P}$为零一损失期望风险，概率同时来自样本抽取和$A$的内部随机性。
\end{theorem}

\begin{proof}
取$2n$个DS打散点及其有限伪立方体$V$，输入均匀，从$V$均匀选目标，并为每向量固定一个实现假设。未见坐标即使给其余全标签仍有至少两个等可能答案，平均错误至少$1/2$。至少一半点未见，总平均风险便至少$1/4$。选择达到平均下界的固定目标，再用$\mathbb E R\leq1/8+(7/8)\mathbb P\{R>1/8\}$，得到所列失败概率。
\end{proof}

有限DS的充分性要构造学习器，不能要求全空间UC。Brukhim等先重构有限标签列表，再从列表得到单标签预测，两步合为压缩\scite{brukhim2022multiclass}。下面先明确借用的构造，再推导它的风险保证。

列表预测$\mu$在每个$x$返回有限集合$\mu(x)\subseteq Y$，例如保留$\{\text{正常},\text{广告}\}$两个答案。它覆盖样本意味着每个真实标签都在当地列表中，尚未完成单标签分类。第二步据列表重构一个全样本一致的单标签规则。

\begin{lemma}[有限DS维的压缩构造]
\label{lem:finite-ds-compression}
设$X,Y$为非空集合，$\mathcal H\subseteq Y^X$非空，且
$d=d_{DS}(\mathcal H)<\infty$。存在通用常数$C>0$，使得对每个整数$n\geq1$，
以及每份大小$q\leq n$、可由某个$h\in\mathcal H$正确分类的带标签样本
$S=((x_i,y_i))_{i=1}^q$，均存在如下压缩表示：从$S$中选取一个允许重复、保留次序的
带标签序列，长度至多
\[
 r_n=\min\!\left\{n,\left\lceil
 C(d+1)^{3/2}\log^2\!\bigl(2n(d+1)\bigr)\right\rceil\right\},
\]
并记录原样本大小$q$。对每个$q\leq n$，存在不依赖具体样本的固定重构映射，仅根据
$q$和该压缩序列输出分类器$\rho_q:X\to Y$，且
$\rho_q(x_i)=y_i$对$i=1,\ldots,q$均成立。重构分类器不要求属于$\mathcal H$。
\end{lemma}

这是原文定理36在常数和对数因子上放宽的表述。引理39用有限DS维重构覆盖原样本的列表，引理40用有限Natarajan维在列表内重构一致单标签规则；拼接压缩序列即可依次重构，$d_N\leq d$使第二步也受$d$控制。取$t=\lceil\sqrt{d+1}\rceil$得到所列对数平方长度；上界超过$n$时保存全样本即可。列表部分借用原文的超图定向，本章仅说明该构造怎样进入压缩与风险分析。

接下来限制可能的压缩序列数，再沿第~\ref{chap:binary-learnability}章的压缩分析同时控制全部重构结果。计数依赖序列长度，未依赖标签总数，因此无限标签仍可分析。

\begin{theorem}[有限DS维的期望风险保证]
\label{thm:finite-ds-learning-rates}
设$(X,\mathcal X)$为可测空间，$Y$为离散标签空间，
$\mathcal H\subseteq Y^X$是由可测分类器组成的非空集合，
$d=d_{DS}(\mathcal H)<\infty$，并采用零一损失。存在通用常数$C>0$；
对每个整数$n\geq1$令
\[
 r_n=\min\!\left\{n,\left\lceil
 C(d+1)^{3/2}\log^2\!\bigl(2n(d+1)\bigr)\right\rceil\right\}.
\]
固定满足$r=r_n<n$的$n$。假定每个候选分类器的零一损失函数，以及所用压缩、
重构和有限候选选择映射均可测。预先固定一个$h_0\in\mathcal H$和全部压缩描述的
顺序；对每份$S=((x_i,y_i))_{i=1}^n$，枚举由原样本大小、至多$r$个允许重复且
保留次序的样本索引组成的描述，用固定重构映射产生候选，无效描述统一返回$h_0$。
令$A^{\mathrm{real}}$选择顺序最先且与$S$一致的候选；若不存在则返回$h_0$。
令$A^{\mathrm{agn}}$选择顺序最先且零一经验风险最小的候选。记
\[
 B_n(\delta)=\log\frac{2(n+1)(r+1)n^r}{\delta}.
\]
这两个确定性算法允许输出$\mathcal H$之外的可测分类器。对任意
$\delta\in(0,1)$，以下保证成立：
\begin{enumerate}
\item 对任意$X\times Y$上的联合分布$\mathcal P$，若存在
$h^\star\in\mathcal H$满足$R_{\mathcal P}(h^\star)=0$，则对
$S\sim\mathcal P^n$，以至少$1-\delta$的概率，
 \[
 R_{\mathcal P}(A^{\mathrm{real}}(S))
 \leq\frac{B_n(\delta)}{n-r}.
 \]
\item 对任意$X\times Y$上的联合分布$\mathcal P$和
$S\sim\mathcal P^n$，以至少$1-\delta$的概率，
 \begin{align*}
 R_{\mathcal P}(A^{\mathrm{agn}}(S))
 &\leq\inf_{h\in\mathcal H}R_{\mathcal P}(h)+\frac rn\\
 &\quad+\sqrt{\frac{B_n(\delta)}{2(n-r)}}
       +\sqrt{\frac{\log(2/\delta)}{2n}}.
 \end{align*}
\end{enumerate}
其中$R_{\mathcal P}$为零一损失期望风险，两项概率都只来自独立同分布样本。
\end{theorem}

\begin{proof}[证明（使用引理~\ref{lem:finite-ds-compression}）]
有序序列允许重复，长度至多$r$时至多$(r+1)n^r$种索引；再计$n+1$种原样本大小，描述总数至多$M=(n+1)(r+1)n^r$。这些索引在看数据前确定，未预先列举无限标签。

先于数据固定$h_0\in\mathcal H$和描述顺序。描述包含原样本大小、序列长度与可重复索引；重构补为总映射，无效描述返回$h_0$，每份样本便有至多$M$个有序候选。可实现算法取首个全一致者，压缩引理保证它存在；不可知算法取首个经验最小者。题设的可测映射使这两种有限选择也可测。

固定索引描述，设实际读取$j\leq r$个不同观测。给定这些值后，候选已确定，其余$n-j$个观测仍iid。风险大于$u$的候选全分对余下者的概率至多$(1-u)^{n-j}\leq e^{-(n-r)u}$。对$M$种描述作并集界，就覆盖首个全一致输出；取$u=B_n(\delta)/(n-r)$得到可实现界，$u\geq1$时则用风险不超过1的平凡保证。

不可知算法在有限重构结果中选首个经验最小者。对固定描述，以$\hat R_{\mathrm{rest}}$记未读取观测的损失平均。条件独立性与单侧Hoeffding，再跨$M$个描述求和，使下式以至少$1-\delta/2$的概率同时成立：
\[
 R_{\mathcal P}(h)\leq\hat R_{\mathrm{rest}}(h)
       +\sqrt{\frac{\log(2M/\delta)}{2(n-r)}}.
\]
删去$j$条观测后平均损失最多增加$j/n$。把全样本均值写成剩余与已读两部分的加权平均，就得$\hat R_{\mathrm{rest}}(h)\leq\hat R_S(h)+j/n\leq\hat R_S(h)+r/n$，可以代入上式。

任意预先固定$g\in\mathcal H$，把它分对的样本组成可实现子样本，压缩引理提供在这些观测上全对的重构者；其经验风险至多$\hat R_S(g)$。子样本为空时，$\hat R_S(g)=1$，任意候选都不更大。经验最小重构于是满足$\hat R_S(A^{\mathrm{agn}}(S))\leq\hat R_S(g)$。对固定$g$再用Hoeffding，以至少$1-\delta/2$的概率有
\[
 \hat R_S(g)\leq R_{\mathcal P}(g)
       +\sqrt{\frac{\log(2/\delta)}{2n}}.
\]
合并两个半失败预算即得到相对该$g$的风险界。任取正的$\tau$，在抽样前选风险距类内下确界不足$\tau$的$g$，得到超过所列上界加$\tau$的概率至多$\delta$；令$\tau\downarrow0$，由事件单调收敛得到第二结论，无需所有参照同时估准。
\end{proof}

固定$d$时$r_n=O(\log^2n)$，所以$r_n/n\to0$，两种风险界分别趋于零与类内最优基准。略去$n,d$的对数因子，可实现风险为$\widetilde O((d^{3/2}+\log(1/\delta))/n)$，不可知超额为$\widetilde O(\sqrt{(d^{3/2}+\log(1/\delta))/n})$。这些量说明统计保证，没有保证压缩搜索可计算或多项式时间完成。

无限DS的必要性与有限DS的构造合起来，得到多类基本定理。前章VC统一描述的性质现在分为图维和DS维两组；有限标签时，两组有限性再度等价，一般$Y$时则须区分。

\begin{theorem}[多分类学习基本定理]
\label{thm:ds-characterization}
设$(X,\mathcal X)$为可测空间，$Y$为取离散$\sigma$代数的非空标签空间，
$\mathcal H\subseteq Y^X$是由可测分类器组成的非空集合，并采用零一损失。假定
\[
 \mathcal L_{\mathcal H}
 =\{(x,y)\mapsto\mathbb I\{h(x)\neq y\}:h\in\mathcal H\}
\]
中的函数均可测，且存在可数子类$\mathcal L_0\subseteq\mathcal L_{\mathcal H}$，
使每个损失函数都是$\mathcal L_0$中某个函数序列的逐点极限。假定有限DS压缩构造中的
压缩、重构和有限候选选择映射均可测。学习算法均为可测映射，允许输出$\mathcal H$
之外的可测预测器；训练样本均为$S\sim\mathcal P^n$，成功概率同时计入样本抽取和
算法的内部随机性。在这一分布无关设定中，以下命题等价：
\begin{enumerate}
 \item $d_{DS}(\mathcal H)<\infty$；
 \item $\mathcal H$可实现PAC可学习；
 \item $\mathcal H$不可知PAC可学习。
\end{enumerate}
此外，$d_G(\mathcal H)<\infty$当且仅当零一损失满足分布无关的一致收敛；
有限图维或一致收敛中的任一条件均蕴含前三条，但一般不存在反向蕴含。
\end{theorem}

\begin{proof}[证明（使用有限DS维的学习器构造结果）]
DS有限时，当样本量足够大，定理~\ref{thm:finite-ds-learning-rates}的两项风险界可分别满足任意给定的误差容限和失败概率要求，因而给出两种分布无关PAC保证。不可知包含可实现，而无限DS已排可实现，故三项等价。图维/UC等价见定理~\ref{thm:graph-uniform-convergence}，学习不推UC由定理~\ref{thm:multiclass-erm-separation}的集合标签例证明。充分性仍依赖已引用的外部学习构造。
\end{proof}

一般多类有两项刻画：图维有限等价于零一UC，DS维有限等价于PAC学习器存在。前者推后者，逆向失败。有限标签使$d_N,d_{DS},d_G$有限性、UC及两种PAC全部等价；二类时三维都回到VC。有限标签的学习等价见定理~\ref{thm:finite-label-learnability}（维度关系可参考\BookTheoremRef{thm:multiclass-dimension-comparison}{第一卷“计数与计算”章的多分类维度比较定理}）。

定性维度有限没有确定最优样本量或高效实现，不能在二类样本式中只换VC为DS便得新公式。

\section{可计算PAC学习}
\label{sec:computable-multiclass-learning}

\subsection{统计保证与程序实现}

沿用第~\ref{sec:computable-binary-learning}节的三层区别：统计学习、训练预测总停机、多项式资源。现在标签程序返回$Y$中的值，仍须停机，风险成功标准保持前述PAC要求。

取$X=\mathbb N$，沿用定义~\ref{def:binary-computable-pac}的程序要求，但须先给标签和假设有效表示：标签编码可计算、相等可判定，有限标签全集可枚举，允许假设及输出由预定可递归枚举表示提供。不可知学习器对每份有限样本停机，返回一个对每输入停机输出$Y$标签的程序，并对所有允许来源和所有足量$n$满足
\[
 \mathbb P_{S\sim\mathcal P^n}\!\left\{
 R_{\mathcal P}(A(S))\leq\inf_{h\in\mathcal H}R_{\mathcal P}(h)+\varepsilon
 \right\}\geq1-\delta
\]
。输出在$\mathcal H$内是恰当，允许其他可计算规则是非恰当。再要求可计算充分样本上界得到strong CPAC，与前章同义，未要求多项式时间。

\begin{theorem}[多分类可计算ERM的充分条件]
\label{thm:multiclass-computable-erm}
在离散输入空间$X=\mathbb N$和具有可计算编码、可判定相等关系的非空标签空间$Y$上
采用零一损失。设非空假设空间$\mathcal H\subseteq Y^X$采用固定有效表示，
并存在确定性程序$A$：它对每份有限样本$S=((x_i,y_i))_{i=1}^n$都停机，输出
$A(S)\in\mathcal H$的程序编码；该输出程序对每个$x\in\mathbb N$都停机，且
\[
 \hat R_S(A(S))
 =\min_{h\in\mathcal H}\frac1n\sum_{i=1}^n\mathbb I\{h(x_i)\neq y_i\}.
\]
若$d_G(\mathcal H)<\infty$，则存在可计算的样本量上界$m(\varepsilon,\delta)$，
使任意$\varepsilon,\delta\in(0,1)$、任意$\mathbb N\times Y$上的联合分布
$\mathcal P$、任意$n\geq m(\varepsilon,\delta)$及
$S\sim\mathcal P^n$均满足
\[
 \mathbb P_{S\sim\mathcal P^n}\left\{
 R_{\mathcal P}(A(S))
 \leq\inf_{h\in\mathcal H}R_{\mathcal P}(h)+\varepsilon
 \right\}\geq1-\delta.
\]
因此，$\mathcal H$恰当地不可知CPAC可学习，且有可计算的样本量上界；概率仅来自
独立同分布样本。若$2\leq|Y|<\infty$且标签全集可有效枚举，则条件
$d_N(\mathcal H)<\infty$可以替代$d_G(\mathcal H)<\infty$。
\end{theorem}

\begin{proof}
定理~\ref{thm:graph-uniform-convergence}将图维等同于损失VC维，因此给出显式UC样本界。取偏差$\varepsilon/2$，两次比较使ERM超额至多$\varepsilon$；训练和预测停机来自题设。固定空间的有限维上界可写入样本量程序。有限标签时，Natarajan维有限蕴含图维有限，所以可替条件（可参考\BookTheoremRef{thm:multiclass-dimension-comparison}{第一卷“计数与计算”章的多分类维度比较定理}）。
\end{proof}

这条充分路线与前章可计算ERM对应：二类用VC，多类用零一损失的图维，有限标签又可用Natarajan。维度有限若没有训练与预测程序，不能仅凭统计定理声称可计算。

\subsection{可计算Natarajan维}

Natarajan有限保证足量输入的固定标签对中有一种二元预测缺失；程序却需要实际找出它。\term{可计算Natarajan维}（computable Natarajan dimension）在这一组合存在性上增加总可执行的见证。

\begin{definition}[可计算Natarajan维]
\label{def:computable-natarajan-dimension}
若存在总可计算程序$W$，对任意$d+1$个互异输入$x_i$及标签对$(a_i,b_i)$、
$a_i\neq b_i$，都输出$s\in\{0,1\}^{d+1}$，使得
\[
 \nexists h\in\mathcal H\quad
 \forall i:\ h(x_i)=\begin{cases}a_i,&s_i=0,\\b_i,&s_i=1,\end{cases}
\]
则称$d$为可计算Natarajan维的一个上界。最小的这种非负整数记为
$c_N(\mathcal H)$；不存在时记为$\infty$。
\end{definition}

普通维度保证缺失存在，可计算维度要求同一程序根据输入与标签对返回缺失者。因此$d_N\leq c_N$；二类时该见证就是前章的有效VC见证。

充分性借用Gourdeau等的有效扩展：有限标签且$c_N(\mathcal H)=d$时，可构造包含$\mathcal H$的$\mathcal H'$，使$d_N(\mathcal H')\leq d+1$，并能有效求出有限输入上的全部可实现预测及实现假设\scite{gourdeau2025multiclass}。下面把该外部构造与后续风险推导分开。

\begin{theorem}[有限标签的非恰当CPAC刻画]
\label{thm:computable-multiclass-characterization}
设$X=\mathbb N$、$2\leq|Y|<\infty$，二者取离散$\sigma$代数，并采用零一损失。
假定标签具有可计算编码、可判定的相等关系，且标签全集可有效枚举；
非空假设空间$\mathcal H\subseteq Y^X$及允许的输出程序采用预先规定的可递归枚举表示。
把$c_N(\mathcal H)$定义为满足以下条件的最小非负整数$d$：存在总可计算程序$W$，
对任意$d+1$个互异输入$x_i$及任意标签对$(a_i,b_i)$、$a_i\neq b_i$，输出
$s\in\{0,1\}^{d+1}$，使不存在$h\in\mathcal H$同时满足
\[
 h(x_i)=\begin{cases}a_i,&s_i=0,\\b_i,&s_i=1\end{cases}
 \qquad(i=1,\ldots,d+1);
\]
若不存在这样的$d$，则令$c_N(\mathcal H)=\infty$。学习器取确定性程序，对每份有限
样本都停机，并可输出有效构造的可递归枚举扩展类中的总可计算假设；PAC成功概率仅来自
$S\sim\mathcal P^n$。在这一表示模型中，$\mathcal H$非恰当地不可知CPAC可学习，
当且仅当$c_N(\mathcal H)<\infty$。
\end{theorem}
\begin{proof}[证明（充分性使用有效扩展引理）]
必要性使用有效的未见输入论证。设CPAC程序$A$存在，固定使精度与失败概率均为$1/16$的保证成立的样本量$n$。给$2n$个输入及各两标签，有限枚举二元目标，再对每目标枚举这组输入上全部长度$n$的有序样本，训练并计算输出在$2n$点的错误比例。全部训练与预测都停机，有限平均因而可计算。

未见输入的随机二元目标分析给出至少一个候选的平均风险不低于$1/4$；有限枚举可找出它并返回选择向量。若它由$\mathcal H$实现，双容限$1/16$的PAC保证反使均值至多$1/16+1/16<1/4$，矛盾，故它是缺失预测。见证规模为$2n$，得到$c_N\leq2n-1$。$n$针对该空间固定写入程序，未要求从任意学习器描述自动算出。

充分性在外部扩展$\mathcal H'$上运行：枚举有限样本的可实现预测，选错误最少者并返回其可计算实现，构成可计算ERM。有限标签及$d_N(\mathcal H')$有限给不可知PAC；又因$\mathcal H\subseteq\mathcal H'$，扩展类风险下确界不大于原类，所以同一输出也满足原类超额要求。输出可能在$\mathcal H$外，故为非恰当学习。
\end{proof}

这里选确定性学习程序，成功概率只含样本。若模型加入随机位，须规定它们的使用，不能直接枚举所有随机运行。必要性与充分性还依赖允许假设和输出程序的有效表示；仅要求输出某个总可计算函数而不定表示，不属于本章刻画范围。

必要方向没有用标签有限，因此相同有效表示模型的无限标签CPAC也须$c_N<\infty$。充分方向的有效扩展依赖有限标签预测枚举，不能沿用到无限$Y$。前述统计反例只算出有限$d_N$，没有算$c_N$，故它不足以判定无限标签的有限$c_N$是否充分，该问题需要另一项定理。

DS有限不自动给伪立方体搜索、列表构造与单标签归约的执行程序，所以统计刻画不能直接命名为CPAC。具体空间仍可走另一充分路线：构造可计算ERM，再证零一UC；有限标签的有限Natarajan或任意标签的有限图维都能支持后一步。

\subsection{高效学习与优化误差}

停机之外还要给样本、训练、预测及输出长度的资源界。多类标签编码也须计费：当$K\geq2$时，类别编号长约$\log K$；单标签任务只需固定常量描述，集合等结构标签则按其实际描述长计算。

标签数不影响近似ERM的三次比较。若输出$A(S)\in\mathcal H$且
\[
 \hat R_S(A(S))\leq\inf_{h\in\mathcal H}\hat R_S(h)+\alpha,
\]
全部风险偏差又至多$\eta$，则
\[
 R_{\mathcal P}(A(S))\leq\inf_{h\in\mathcal H}R_{\mathcal P}(h)+2\eta+\alpha.
\]
按定理~\ref{thm:approximate-erm-risk}，$2\eta$来自两次抽样均值替换，$\alpha$来自经验目标求解差。若用规模参数的多项式样本使$\eta\leq\varepsilon/4$，以多项式训练时间取得$\alpha\leq\varepsilon/2$，并使输出长和预测时间也为多项式，就得到高效不可知PAC的充分路线。

\subsection{从枚举到有效算法}

前章合取用结构搜索省去指数枚举，多类也能这样做。取$X=\{1,\ldots,p\},Y=\{1,\ldots,K\}$及全部函数$\mathcal H=Y^X$，候选虽有$K^p$个，经验目标却可按输入分别求解。

对每个$x$计$N_{x,y}=|\{i:x_i=x,y_i=y\}|$，选计数最多的标签，平局取最小编号，未见输入取1。正确数对$x$可加，不同输入的选择不互相约束，故逐行最优就是全表经验最优。

\begin{theorem}[有限输入上的多分类表学习]
\label{thm:multiclass-table-erm}
设$p,K,n\geq1$，$X=\{1,\ldots,p\}$、$Y=\{1,\ldots,K\}$均取离散
$\sigma$代数，$\mathcal H=Y^X$，并采用零一损失。令$\mathcal P$为
$X\times Y$上的任意联合分布，$S=((x_i,y_i))_{i=1}^n\sim\mathcal P^n$。
确定性学习算法对每个$x\in X,y\in Y$计算
$N_{x,y}=|\{i:x_i=x,\ y_i=y\}|$，在每个$x$处选择计数最大的标签，平局时选择编号
最小者，未出现的输入预测为1。该算法是$Y^X$上的恰当ERM，需要$O(n+pK)$次计数和
比较操作，输出一张含$p$个标签的预测表。对任意
$\varepsilon,\delta\in(0,1)$，若
\[
 n\geq\frac{2}{\varepsilon^2}
       \left(p\log K+\log\frac2\delta\right),
\]
则以关于$S\sim\mathcal P^n$至少$1-\delta$的概率，其零一损失期望风险不超过
$\inf_{h\in\mathcal H}R_{\mathcal P}(h)+\varepsilon$。
\end{theorem}

\begin{proof}
$h$分对样本数为$\sum_{x=1}^pN_{x,h(x)}$，逐行选最大计数即使总和最大，也就使经验风险最小。初始化$pK$计数器、扫描$n$观测、再逐行比较，操作数为$O(n+pK)$。有限类大小$K^p$与Hoeffding并集使偏差超过$\varepsilon/2$的概率至多$2K^pe^{-n\varepsilon^2/2}$，令其至多$\delta$再用ERM两次比较，得到风险保证。
\end{proof}

操作按$p,K$计费，计数器至多$O(\log(n+1))$位，所以比特时间也为这些参数的多项式。若巨大$K$只由短程序表示，$O(K)$却不等于其描述长的多项式。当$K\geq2$时，输出表长$O(p\log K)$也进入资源预算；$K=1$时，一张含$p$个固定标签的表可用$O(p)$位表示，也可用常量规则表示。

\section{本章小结}
\label{sec:multiclass-learning-summary}

多类让同一输入拥有多个替代答案，改变了前章的容量刻画。编码归约仍须分析子类复杂度、位错累积、无效码解法及非恰当输出。在有限标签下，Natarajan限制预测数，计数、对称化和ERM比较因而沿用；在本章可测条件下，其有限性、分布无关PAC与UC等价。

无限标签分开全类估计与算法存在。图维是零一损失类VC维，判定共同风险估计；DS维判定统计学习，学习器无须全类估准。Natarajan有限只保留必要性。图维有限足以学习，学习却不推出它有限，这解释了二类等价不能原样外推。

经验最小值相同，选择规则仍可能使泛化不同，ERM目标本身没有说明全部学习行为。训练与预测要求总停机后，还需有效组合见证；有限标签的可计算Natarajan提供刻画，却未给多项式时间。下一章换为一般损失，并用稳定性继续考察UC之外的保证。

\paragraph{延伸阅读}
Daniely等的研究可用于对照成功学习器存在与任意ERM成功，特别核对多个经验最优解的选择怎样改变样本量，以及图维在其中的作用\scite{daniely2015multiclass}。

Brukhim等建立DS学习刻画。本章将列表、压缩与单标签输出连接风险，借用的组合构造则需继续读原文的伪立方体和列表分析；有限标签Natarajan证明不能替代这些结果\scite{brukhim2022multiclass}。

Gourdeau、Lechner与Urner的工作提供对应第~\ref{sec:computable-binary-learning}节有效VC的多类计算视角。阅读时须核标签编码、输出求值以及恰当/非恰当范围，才能判断同一统计空间的实现要求\scite{gourdeau2025multiclass}。
